\section{Evaluation}
\label{sec:evaluation}

This chapter presents a comprehensive evaluation of the adapted inconsistency measures and their implementation, assessing both computational feasibility and practical diagnostic value. The evaluation demonstrates the effectiveness of the approach through systematic experimentation on benchmark planning problems and analysis of the diagnostic insights provided.

\subsection{Experimental Framework and Methodology}

This section establishes the experimental framework used for evaluation, including the selection of benchmark planning problems from established collections, the design of test scenarios covering different types of unsolvable planning instances, and the methodology for measuring both computational performance and diagnostic utility.

\subsection{Benchmark Problem Selection}

This section describes the systematic selection of planning problems from standard benchmarks, ensuring coverage of different domain characteristics, problem sizes, and types of unsolvability. The selection includes problems from International Planning Competition datasets and specifically constructed unsolvable instances representing various conflict patterns.

\subsection{Computational Feasibility Assessment}

This section evaluates the computational feasibility and scalability of the implemented measures on problems of varying complexity. Performance metrics include computation time, memory usage, and scalability limits, with analysis of how problem characteristics affect computational requirements.

\subsection{Diagnostic Utility Analysis}

This section analyzes the diagnostic utility of the implemented measures for distinguishing different types of unsolvable planning problems. This includes evaluation of how effectively the measures identify conflict sources, distinguish between different types of planning failures, and provide meaningful quantitative assessments.

\subsection{Comparative Analysis}

This section presents comparative analysis between different adapted measures, examining their relative strengths and weaknesses for various types of planning problems. The comparison considers both computational efficiency and diagnostic discriminative power.

\subsection{Case Studies}

Detailed case studies demonstrate the application of the measures to representative unsolvable planning problems, showing how the quantitative diagnostics help identify conflict sources and suggest potential repair strategies. The case studies cover different planning domains and conflict types.

\subsection{Validation and Verification}

This section addresses the validation of results through comparison with existing diagnostic approaches where available, verification of measure properties through constructed test cases, and assessment of the consistency and reliability of the diagnostic outputs.

\subsection{Limitations and Scope}

This section discusses the limitations of the current approach, including computational constraints, scope restrictions, and cases where the measures may provide limited diagnostic value. This analysis informs future research directions and practical deployment considerations.
